\SixthArticle{16}
We shall apply this investigation more generally to an infinite series $M,\allowbreak M',\allowbreak M'',\allowbreak M'''$ etc. formed so that the quotients $\frac{M'}{M},\frac{M''}{M'},\frac{M'''}{M''}$ etc. are respectively the values of the fraction
\[
\frac{t^{\lambda}+At^{\lambda-1}+Bt^{\lambda-2}+Ct^{\lambda-3}+\mathrm{etc.}}
{t^{\lambda}+at^{\lambda-1}+bt^{\lambda-2}+ct^{\lambda-3}+\mathrm{etc.}}
\]
for $t=m$, $t=m+1$, $t=m+2$ etc. For brevity we shall denote the numerator of this fraction by $P$ and its denominator by $p$. We shall also suppose that $P,p$ are not identical, or that the differences $A-a,B-b,C-c$ etc. do not all vanish together.

I. Whenever the first nonvanishing one of the differences $A-a,B-b,C-c$ etc. is positive, a bound $l$ can be assigned such that, as soon as the value of $t$ exceeds it, the values of the functions $P$ and $p$ will certainly always be positive, with $P>p$. This evidently occurs if $l$ is taken to be the largest real root of the equation $p(P-p)=0$; if that equation has no real roots at all, the property holds for all values of $t$. Thus, in the series $\frac{M'}{M},\frac{M''}{M'},\frac{M'''}{M''}$ etc., all the terms will be positive and greater than unity at least after a certain interval, if not from the start. Hence, if none becomes either $=0$ or infinitely large, it is clear that
\SourceLeaf{149}
\textit{all the terms of the series $M,M',M'',M'''$ etc. will have the same sign and steadily increase, if not from the start, at least after a certain interval.}

In the same way, if the first nonvanishing one of the differences $A-a,B-b,C-c$ etc. is negative, all the terms of the series $M,M',M'',M'''$ etc. will have the same sign and steadily decrease, if not from the start, at least after a certain interval.

II. If the coefficients $A,a$ are unequal, the terms of the series $M,\allowbreak M',\allowbreak M'',\allowbreak M'''$ etc. will either grow without bound or diminish towards zero, according as the difference $A-a$ is positive or negative. We prove this as follows. If $A-a$ is a positive quantity, take an integer $h$ such that $h(A-a)>1$, and set $\frac{M^h}{m}=N$, $\frac{M^{\prime h}}{m+1}=N'$, $\frac{M^{\prime\prime h}}{m+2}=N''$, $\frac{M^{\prime\prime\prime h}}{m+3}=N'''$ etc., and also $tP^h=Q$, $(t+1)p^h=q$. Then $\frac{N'}{N},\frac{N''}{N'},\frac{N'''}{N''}$ etc. are evidently the values of the fraction $\frac{Q}{q}$ on setting $t=m$, $t=m+1$, $t=m+2$ etc., while $Q,q$ are algebraic functions of the following form:
\[
\begin{array}{l}
Q=t^{\lambda h+1}+hAt^{\lambda h}+\mathrm{etc.}\\
q=t^{\lambda h+1}+(ha+1)t^{\lambda h}+\mathrm{etc.}
\end{array}
\]
Since, by hypothesis, the difference $hA-(ha+1)$ is a positive quantity, the terms of the series $N,N',N'',N'''$ etc. will steadily increase, if not from the start, at least after a certain interval (by \textsc{i}). Hence the terms of the series $mN,(m+1)N',(m+2)N''$, $(m+3)N'''$ etc. must grow without bound, and so must the terms of the series $M,M',M'',M'''$ etc., whose powers with exponent $h$ equal those terms. Q.~E.~P.

If $A-a$ is a negative quantity, an integer $h$ must be taken so that $h(a-A)$ becomes greater than $1$. Then, by similar reasoning, the terms of the series
\[
mM^h,\quad (m+1)M^{\prime h},\quad (m+2)M^{\prime\prime h},\quad (m+3)M^{\prime\prime\prime h}\quad\mathrm{etc.}
\]
will steadily decrease after a certain interval. Thus the terms of the series $M^h$, $M^{\prime h},M^{\prime\prime h}$ etc., and therefore also the terms of $M,M',M'',M'''$ etc., must diminish towards zero. Q.~E.~S.

III. If, however, the first coefficients $A,a$ are equal, the terms of the series $M,\allowbreak M',\allowbreak M'',\allowbreak M'''$ etc. will steadily converge towards a finite limit. We prove this as follows. First we suppose that the terms of the series steadily increase after a certain interval, or that the first nonvanishing one of the differences $B-b,C-c$ etc.
\SourceLeaf{150}
is positive. Let $h$ be an integer such that $h+b-B$ becomes a positive quantity, and set
\[
M\left(\frac{m}{m-1}\right)^h=N,\quad
M'\left(\frac{m+1}{m}\right)^h=N',\quad
M''\left(\frac{m+2}{m+1}\right)^h=N''\quad\mathrm{etc.}
\]
and $(tt-1)^hP=Q,\ t^{2h}p=q$, so that $\frac{N'}{N},\frac{N''}{N'}$ etc. are the values of the fraction $\frac{Q}{q}$ on setting $t=m,\ t=m+1$ etc. Since we then have
\[
\begin{array}{l}
Q=t^{\lambda+2h}+At^{\lambda+2h-1}+(B-h)t^{\lambda+2h-2}\quad\mathrm{etc.}\\
q=t^{\lambda+2h}+At^{\lambda+2h-1}+bt^{\lambda+2h-2}\quad\mathrm{etc.}
\end{array}
\]
and, by hypothesis, $B-h-b$ is a negative quantity, the terms of the series $N,N',N'',N'''$ etc. will steadily decrease at least after a certain interval. Thus the terms of the series $M,M',M'',M'''$ etc., which are always respectively smaller than those terms, can only converge towards a finite limit as they steadily increase. Q.~E.~P.

If the terms of the series $M,M',M'',M'''$ etc. steadily decrease after a certain interval, an integer must be taken for $h$ such that $h+B-b$ is a positive quantity. Entirely similar reasoning will show that the terms of the series
\[
M\left(\frac{m-1}{m}\right)^h,\quad
M'\left(\frac{m}{m+1}\right)^h,\quad
M''\left(\frac{m+1}{m+2}\right)^h\quad\mathrm{etc.}
\]
steadily increase after a certain interval. Hence the terms of the series $M,M',M''$ etc., which are always respectively greater than those terms, can only decrease towards a finite limit as they steadily decrease. Q.~E.~S.

IV. Finally, as regards the \textit{sum} of the series whose terms are $M,M',$ $M'',M'''$ etc., in the case where they diminish towards zero, we first suppose that $A-a$ lies between $0$ and $-1$, or that $A+1-a$ is either a positive quantity or $=0$. Let $h$ be a positive integer, arbitrary when $A+1-a$ is a positive quantity, or such that it makes the quantity $h+m+A+B-b$ positive when $A+1-a=0$. Then
\[
\begin{array}{l}
P(t-(m+h-1))=t^{\lambda+1}+(A+1-m-h)t^\lambda\\\qquad{}+(B-A(m+h-1))t^{\lambda-1}\quad\mathrm{etc.}\\
p(t-(m+h))=t^{\lambda+1}+(a-m-h)t^\lambda\\\qquad{}+(b-a(m+h))t^{\lambda-1}\quad\mathrm{etc.}
\end{array}
\]
where either $A+1-m-h-(a-m-h)$ will be a positive quantity, or, if it becomes $=0$, at least $B-A(m+h-1)-(b-a(m+h))$ will be positive. Hence (by \textsc{i}) for the quantity $t$ a value $l$ can be assigned such that, as soon as it exceeds that value, the values of the fraction $\frac{P(t-(m+h-1))}{p(t-(m+h))}$ will always become positive and greater than unity. Let $n$ be an integer
\SourceLeaf{151}
greater than $l$ and also greater than $h$, and let the terms of the series $M,M',M'',M'''$ etc. corresponding to the values $t=m+n,\ t=m+n+1,\ t=m+n+2$ etc. be $N,N',N'',N'''$ etc. Then
\[
\frac{(n+1-h)N'}{(n-h)N},\quad
\frac{(n+2-h)N''}{(n+1-h)N'},\quad
\frac{(n+3-h)N'''}{(n+2-h)N''}\quad\mathrm{etc.}
\]
are positive quantities greater than unity, whence
\[
N'>\frac{(n-h)N}{n+1-h},\quad
N''>\frac{(n-h)N}{n+2-h},\quad
N'''>\frac{(n-h)N}{n+3-h}\quad\mathrm{etc.}
\]
and hence the sum of the series $N+N'+N''+N'''+$ etc. is greater than the sum of the series
\[
(n-h)N\left(\frac{1}{n-h}+\frac{1}{n+1-h}+\frac{1}{n+2-h}+\frac{1}{n+3-h}+\mathrm{etc.}\right)
\]
however many terms are collected. But the latter series exceeds every bound as the number of terms grows to infinity, since the sum of the series $1+\frac12+\frac13+\frac14+$ etc., which is known to be infinite, remains infinite even if the initial terms $1+\frac12+\frac13+\mathrm{etc.}$ $+\frac{1}{n-1-h}$ are removed. Thus the sum of the series $N+N'+N''+N'''+$ etc., and therefore also the sum of $M+M'+M''+M'''+$ etc., of which it is a part, grows without bound.

V. Whenever $A-a$ is a negative quantity greater than unity in absolute value, the sum of the series $M+M'+M''+M'''+$ etc. continued to infinity will certainly be finite. For let $h$ be a positive quantity smaller than $a-A-1$. Similar reasoning will show that a value $l$ of the quantity $t$ can be assigned beyond which the fraction $\frac{Pt}{p(t-h-1)}$ always takes positive values smaller than unity. If $n$ is now taken to be an integer greater than $l,m,h+1$, and the terms of the series $M,M',M'',M'''$ etc. corresponding to the values $t=n,\ t=n+1,\ t=n+2$ etc. are denoted by $N,N',N''$ etc., then
\[
N'<\frac{n-h-1}{n}\cdot N,\quad
N''<\frac{(n-h-1)(n-h)}{n(n+1)}\cdot N'\quad\mathrm{etc.}
\]
and hence the sum of the series $N+N'+N''+$ etc., however many terms are collected, is smaller than the product of $N$ and the sum of the same number of terms of the series
\[
1+\frac{n-h-1}{n}+\frac{(n-h-1)(n-h)}{n(n+1)}+
\frac{(n-h-1)(n-h)(n-h+1)}{n(n+1)(n+2)}\quad\mathrm{etc.}
\]
Its sum can readily be given for any number of terms, namely
\[
\begin{array}{lr}
\mathrm{first\ term}&=\displaystyle\frac{n-1}{h}-\frac{n-h-1}{h}
\end{array}
\]
\SourceLeaf{152}
\[
\begin{array}{@{}l@{\;}l@{}}
\text{sum of two terms}&=\displaystyle\frac{n-1}{h}-\frac{(n-h-1)(n-h)}{hn}\\[7pt]
\text{sum of three terms}&=\displaystyle\frac{n-1}{h}-\frac{(n-h-1)(n-h)(n-h+1)}{hn(n+1)}\quad\text{etc.}
\end{array}
\]
and since the second part (by II) forms a series diminishing towards zero, that sum continued to infinity must be set $=\frac{n-1}{h}$. Thus $N+N'+N''$ etc. continued to infinity will always remain smaller than $\frac{N(n-1)}{h}$, and therefore $M+M'+M''$ etc. will certainly converge to a finite sum. Q.~E.~D.

VI. To apply what we have proved in general about the series $M,M',M''$ etc. to the coefficients of the powers $x^m,x^{m+1},x^{m+2}$ etc. in the series $F(\alpha,\beta,\gamma,x)$, it will be necessary to set $\lambda=2$, $A=\alpha+\beta$, $B=\alpha\beta$, $a=\gamma+1$, $b=\gamma$. The five assertions proposed in the preceding article then follow at once.

